% !PRINTED 266
\weberpage{266}
%<<CONT
Paragraphen durch die Vertauschung $\Bigl(u,\dfrac{\sqrt2}{u}\Bigr)$ die
Gesamtheit der Wurzeln $v$ in
\[
\Bigl(\frac{2}{p}\Bigr)\frac{\sqrt2}{v}
\]
\"uber. Hieraus folgt, da\ss\ f\"ur $u=0$ die s\"amtlichen Wurzeln $v$
in Null \"ubergehen, also keine von ihnen unendlich wird, woraus
zu schlie\ss en ist, da\ss\ nicht nur die $U_1, U_2, \ldots U_{p+1}$
\gesperrt{ganze} Funktionen von $u$ sind, sondern da\ss\ auch jede
von ihnen den Faktor $u$ enthalten mu\ss.

Wir schlie\ss en nun zun\"achst, genau wie bei der Invariantengleichung
(\S\ 69, 2., 3. und 6.), da\ss
\weq{5}{\Phi_p(v,u)=\Phi_p(u,v),}
und da\ss
\weq{6}{\Phi_p(v,u)=(v^p-u)(v-u^p)+p\sum_{1,p-1}^{h,k}c_{h,k}u^hv^k,}
worin $c_{h,k}=c_{k,h}$ \gesperrt{ganze Zahlen} sind.
Wir wollen diese Funktion in der Weise darstellen
\weq{7}{\Phi_p(v,u)=v^{p+1}+u^{p+1}+\sum_{1,p}^{h,k}a_{h,k}u^hv^k,}
worin also $a_{h,k}=a_{k,h}$ ebenfalls ganze Zahlen sind, auf deren
Teilbarkeit durch $p$ es nun weiter nicht ankommt, und es ist
insbesondere $a_{p,p}=-1$. Wenn wir in der Gleichung
\[
\Phi_p[f(p\,\omega), f(\omega)]=0
\]
$\omega$ durch $\omega+2$ ersetzen, so ergibt sich wegen der
Irreducibilit\"at auf Grund der Relation
\[
f(\omega+2)=e^{-\frac{\pi i}{12}}f(\omega),
\]
da\ss\ in (7) nicht \gesperrt{alle} Glieder, sondern nur solche vorkommen,
in denen die Exponenten $h, k$ der Kongruenz
\weq{8}{hp+k\equiv p+1\ (\mathrm{mod}\ 24)}
gen\"ugen.

Nun kennen wir noch eine weitere Eigenschaft der Funktionen
$\Phi_p(v,u)$, die sich aus der schon benutzten Vertauschung (9)
des vorigen Paragraphen ergibt, wo $\varrho=1$ zu setzen ist, und
die, wenn wir zur Abk\"urzung
\[
\varepsilon=\Bigl(\frac{2}{p}\Bigr)=(-1)^{\frac{p^2-1}{8}}
\]
%CONT>>
